% The Futhark domain paper of the multi-typed GEP set.
% Build: cd papers && latexmk -pdf futhark_multityped_gep.tex
% Every number in this paper has a source file named beside it; the figure data are in
% futhark_multityped_gep_data/ (each file names the ledger it was taken from).
\documentclass[sigconf,nonacm]{acmart}

\usepackage{booktabs}
\usepackage{tabularx}
\usepackage{pgfplots}
\pgfplotsset{compat=1.18}
\usepackage{xcolor}
\usepackage{listings}

\settopmatter{printacmref=false,printfolios=true}
\setlength{\emergencystretch}{1.5em}
\renewcommand\footnotetextcopyrightpermission[1]{}

\DeclareRobustCommand{\src}[1]{{\footnotesize\path{#1}}}
\newcommand{\ty}[1]{\texttt{#1}}
\newcommand{\D}{\textsf{D}}
\newcommand{\Hh}{\textsf{H}}
\newcolumntype{Y}{>{\raggedright\arraybackslash}X}
\newcommand{\dat}[1]{futhark_multityped_gep_data/#1}

\lstset{basicstyle=\ttfamily\scriptsize,columns=fullflexible,breaklines=true,frame=none,
  xleftmargin=0pt,aboveskip=3pt,belowskip=3pt}

\definecolor{evo}{HTML}{1f5fa8}
\definecolor{rnd}{HTML}{c0392b}

\begin{document}

\title{The Futhark Kingdom: Multi-typed Gene Expression Programming for GPU Array Programs}
\subtitle{A benchmark study, the fitness faults it found, and what the corrected fitness does}

\author{Andrew James Morgan}
\affiliation{\institution{Gamakon Ltd.}\country{United Kingdom}}
\email{andrew@gamakon.ai}

\begin{abstract}
The Futhark kingdom applies Andrew Morgan's multi-typed gene expression
programming (GEP) to Futhark, a purely functional array language that compiles
to GPU code. Its rules of legality are types: a matrix's shape class, an index
that is in range, a variable that is bound, and a bounded amount of work. Every
gene therefore expresses a program the Futhark compiler accepts. The kingdom has
three families, all generated offline from declared interfaces: a matrix
family (11 types, 71 rows), a signature family that builds one kingdom from a
benchmark's typed signature (73 to 832 types), and a pixel family for image
shaders (435 types). Candidates are compiled by a fork of Futhark's WebGPU
compiler that we extended with a Rust and wgpu backend. They run and are timed
on the GPU, and the host scores them with Morgan's hyperspherical fitness
function (HFF).

We report four measured results. (1) On a matrix kernel that computes a product
twice, the search found a form with 9 device allocations instead of 15. The
reference's OpenCL kernel was too large for the device to build, and the found
form ran. (2) In a study on four published benchmarks, evolution was no better
than random search. Both found the two-node reduce and scan exactly; neither
found histogram or pathfinder in 1,400 evaluations. (3) A review of that study
found 22 faults, six of which break evolution: fitness got \emph{worse} as
programs got closer to the published one (Spearman $-0.68$ and $-0.52$). We added
a slope check that every run must pass first, rebuilt the quality inputs and
the fitness cases, and reran. Evolution then reached a lower best HFF than
random search in 4 of 4 pairs, and its median HFF fell in every run. No run
found an exact program. (4) No program found by evolution is faster than the
program it started from. One pathfinder program ran 5\,\% faster than the
published port, but that speed comes from the way the kingdom writes the
algorithm, not from the search. We also fixed the compiler and timing faults
that the study exposed: SetMem in GPU kernels, timing that drifted with the
GPU's clock, and three evaluation paths that disagreed.
\end{abstract}

\keywords{gene expression programming, strongly typed genetic programming, Futhark,
GPU, superoptimisation, fitness landscape, equality saturation}

\maketitle

%======================================================================
\section{Introduction}

Genetic programming over a real programming language has two problems before
it has a search problem. Most random programs are illegal, and the legal ones
must be compiled and run before they can be scored. Andrew Morgan's
multi-typed GEP removes the first problem by construction. A Karva gene has a
head and a tail for each output type, and a slot that needs type $T$ reads the
next row of $T$'s own reading frame. Every expressed node therefore has the type
its parent asks for, so every gene expresses a legal program. The schema,
the master symbol table, the geneframe and the GPU engine are described in
the core paper~\cite{morgan_core}. This paper covers one domain:
Futhark~\cite{henriksen2017futhark}.

Futhark is a good target for three reasons. Its programs are pure functions of
typed arrays, so a candidate's behaviour is fully described by its outputs on a
set of inputs. Its type system already knows about array sizes, so the rules
that make a program legal can be written as types. And its optimising compiler
does fusion, flattening and tiling~\cite{henriksen2019flattening}, so the search
can concentrate on what the compiler cannot see.

The paper follows the work in the order it happened, failures included,
because the failures were the main finding:
\begin{enumerate}
\item a matrix kernel on which the e-graph and the compiler's common
  subexpression elimination together turn a kernel the device cannot build into
  one it can (Section~\ref{sec:a7});
\item a study on four published benchmarks, in which evolution was no better
  than random search (Section~\ref{sec:study});
\item the 22 faults a review then found, chiefly a fitness that rewarded moving
  \emph{away} from the target, and a slope check that now has to pass before any
  run (Section~\ref{sec:faults});
\item reruns on the corrected fitness, in which evolution beat random search in
  4 of 4 pairs but found no exact program (Section~\ref{sec:reruns});
\item the compiler and timing faults the study exposed, and their fixes
  (Section~\ref{sec:fixes}).
\end{enumerate}

\paragraph{Contributions.}
\begin{itemize}
\item \textbf{A Futhark kingdom in which legality is a type.} Shape classes,
  in-range indices, bound variables and a work bound are all types. They are
  generated offline from a declared interface, and witness tests re-check every
  relation from the emitted tables (Section~\ref{sec:kingdom}).
\item \textbf{An evaluation path that scores the code that ships.} Our fork's
  \texttt{rust-wgpu} backend compiles each generation into one Rust module that
  runs inside the evolution process on our own wgpu device. Its outputs are
  bit-identical to \texttt{futhark c} on the integer benchmarks
  (Section~\ref{sec:gpu}).
\item \textbf{A measured account of a fitness that could not guide a search}, and
  a slope check that refuses to start a run on such a fitness
  (Sections~\ref{sec:faults} and~\ref{sec:slope}).
\item \textbf{Controlled comparisons against random search}, run with equal
  numbers of new programs and timed with interleaved pairs and an A/A noise floor
  (Sections~\ref{sec:study} and~\ref{sec:reruns}).
\end{itemize}

%======================================================================
\section{Related work}

\paragraph{Typed genetic programming and GEP.}
Genetic programming evolves program trees~\cite{koza1992gp}. Strongly typed
GP~\cite{montana1995stgp} keeps trees well typed by constraining generation and
variation. GEP~\cite{ferreira2001gep} separates a fixed-length linear genotype
from the tree it expresses. Morgan's multi-typed GEP~\cite{morgan_core} has a
head and tail for every type, so type safety is a property of the layout rather
than of the operators.

\paragraph{GP on GPUs.}
Early GPU GP evaluated a population of programs in parallel on graphics
hardware~\cite{harding2007gpu,chitty2007gpu}. In the Futhark kingdom the GPU
runs the \emph{evolved programs themselves}, compiled by a production compiler.
In one of the two runners it also holds the population (Section~\ref{sec:gpu}).

\paragraph{Superoptimisation.}
STOKE~\cite{schkufza2013stoke} searches for x86 code by stochastic search, and
Souper~\cite{sasnauskas2017souper} synthesises LLVM IR with a solver. We know of
no superoptimiser for Futhark. Our matrix result
(Section~\ref{sec:a7}, and the short paper~\cite{morgan2026a7}) is of this kind,
but on a single kernel.

\paragraph{Equality saturation.}
E-graphs hold many equal forms of a program at once~\cite{tate2009equality,willsey2021egg}.
We use egglog~\cite{zhang2023egglog}, with one sort for each kingdom type and
rules that are identities.

\paragraph{Benchmarks and landscape measures.}
The benchmarks are Futhark ports~\cite{futharkbenchmarks} of
Rodinia~\cite{che2009rodinia} programs and of micro benchmarks in the style of
Parboil~\cite{stratton2012parboil}. Our slope check is a form of fitness-distance
correlation~\cite{jones1995fdc}, with the tree edit distance of Zhang and
Shasha~\cite{zhang1989ted}. Futhark's size
types~\cite{henriksen2021sizedependent} are related to the matrix family's
classes; the classes are generated by our own table generator,
not by Futhark's type system.

%======================================================================
\section{The Futhark kingdom}
\label{sec:kingdom}

\subsection{What is shared, and what is not}
The genotype, the head and tail lengths for a declared program size $N$
($h_T=[T{=}R]+(N-1)\,a_T$), expression, and the contract that every gene
expresses a legal program all come from the core paper~\cite{morgan_core}.
What belongs to this kingdom is its \emph{tables}. They are written once,
offline, by \src{kingdoms/futhark_shapes/gen_tables.py}, from an interface
that is declared in the generator. The engine reads them as an ordinary flat
typed kingdom and never learns that shapes exist. The tables are therefore the
trust boundary, and two checks guard them. The generator's default mode,
\texttt{--check}, compares every committed table byte for byte with what it
would write now, and a Rust test runs that check, so a hand edit fails
\texttt{cargo test}. Then witness tests re-read the committed text and re-check
every shape relation against a separately written model.

The kingdom has three families, all from the same generator
(Table~\ref{tab:families}).

\begin{table}[t]
\caption{The Futhark kingdom's instances, counted from the committed tables
(\src{kingdoms/futhark_shapes/}: \texttt{types.tsv}, \texttt{functions.tsv}, \texttt{egraph.rules}).}
\label{tab:families}
\small
\begin{tabular}{@{}llrrr@{}}
\toprule
family & instance & types & rows & e-graph rules\\
\midrule
matrix    & A7 interface ($n,m,p$)        & 11  & 71    & 54\\
signature & reduce \texttt{[n]i32 -> i32}   & 73  & 572   & 219\\
signature & scan \texttt{[n]i32 -> [n]i32}  & 73  & 590   & 219\\
signature & histogram                     & 607 & 4,389 & 1,041\\
signature & pathfinder                    & 832 & 5,457 & 1,545\\
pixel     & image target                  & 435 & 4,957 & 410\\
pixel     & cuneus                        & 435 & 5,531 & 451\\
pixel     & volumepassage                 & 435 & 5,490 & 451\\
\bottomrule
\end{tabular}
\end{table}

\subsection{The matrix family: a shape class is a type}
The interface declares three dimension roles, $n$, $m$ and $p$. A class
\ty{mat.n.m} is the interface's \texttt{[n][m]f32}. There are nine matrix
classes and two function classes: a row function \texttt{[]f32 -> []f32}, used
by \texttt{map} over rows, and a binary row function, needed because Futhark's
\texttt{map2} is curried. The rows are \texttt{map\_rows}, \texttt{map2\_rows}
and \texttt{transpose} for every class (27 rows); \texttt{matmul} for every
legal triple of classes (27); one input terminal for each class, plus two
\texttt{const} rows (11); and six row functions. The output class is
\ty{mat.n.p}. The program renders as a shape-polymorphic Futhark function over
the roles, and an entry point binds the actual sizes. One gene is therefore a
program for any sizes. At $N=3$ the gene is 322 words long.

Futhark 0.28.1 has no \texttt{matmul}, so the kingdom brings a
shape-polymorphic helper:
\begin{lstlisting}
def matmul [n][m][p] (A: [n][m]f32) (B: [m][p]f32) : [n][p]f32 =
  map (\a -> map (\b -> reduce (+) 0.0 (map2 (*) a b)) (transpose B)) A
\end{lstlisting}
Futhark's own type checker then enforces the inner-dimension relation that the
class table states. The table and the compiler check each other: the
compiler caught the table giving \texttt{map2} the wrong arity
(\src{kingdoms/futhark_shapes/README.md}).

The e-graph has one egglog sort for each class, so a rewrite cannot change a
node's class. It has 54 rewrites, all bit-identical. For each class there are
27 identities: \texttt{map} with the identity row function is the matrix;
\texttt{map2} with the binary identity is its first operand; and
\texttt{transpose} is an involution. There are also 27 transposed-product rules,
$(Y^\top X^\top)^\top = XY$ (\src{kingdoms/futhark_shapes/egraph.rules}). (The
kingdom's README still says 27 rules; the committed table has 54.)

Two costs rank the programs. The \emph{static cost} is a polynomial in the
roles (traffic, allocations and serial depth), generated into
\texttt{costs.tsv} with a witness test of its own
(\src{src/futhark_shapes_cost.rs:263}). The \emph{measured GPU cost} comes from
the compiled GPU program (\texttt{futhark dev --gpu-mem --compile-imp-gpu}). It
records device allocation sites, literal bytes, shared memory, kernels, and
whether the block and register tiler fired
(\src{src/futhark_shapes_gpu.rs:110}).

\subsection{The signature family: a kingdom from a benchmark's signature}
\label{sec:sig}
A signature instance declares a program's typed inputs and output over named
dimensions, for example Rodinia pathfinder's
\texttt{walls: [rows][cols]i32 -> [cols]i32}. The generator derives the rest,
and makes every rule of a legal program a type
(\src{kingdoms/futhark_shapes/gen_tables.py:1684-1811}):
\begin{itemize}
\item \textbf{A class is a type.} \ty{a.rows.cols}, \ty{a.cols}, \ty{ix.cols}
  and \ty{i32} are classes; a row of an \ty{a.rows.cols} is an \ty{a.cols}.
\item \textbf{An index is in range by its type.} \ty{ix.D} has the terminals
  \texttt{0}, $D{-}1$ and the counters of the binders over $D$. Its functions are
  \texttt{prev} ($\max(0,i{-}1)$), \texttt{next} ($\min(D{-}1,i{+}1)$) and
  \texttt{clamp}. Every dimension is at least 1, so every index is legal, and
  neighbour accesses at the edges are clamped by the types.
\item \textbf{A variable is bound before it is read.} A type is a class in a
  context, the stack of binders around the node: \texttt{T-D} (a
  \texttt{tabulate D}), \texttt{L-D-S} (a loop over $D$ carrying state $S$) and
  \texttt{F-D-S} (the same loop from 1). For example,
  \ty{i32:c/\allowbreak F-rows-a.cols/\allowbreak T-cols} is an \texttt{i32} inside a \texttt{tabulate}
  over \texttt{cols}, inside a row loop. Its terminals are the variables of its
  context and the interface's inputs.
\item \textbf{The work is bounded by the types.} Binders, reductions,
  \texttt{map2} and \texttt{replicate} exist only in contexts shallower than
  \texttt{SIG\_DEPTH} $=2$. So a node runs at most $\max(D)^2$ times, and no gene
  can express a runaway kernel.
\item \textbf{Optional families.} An instance can name \texttt{scan},
  \texttt{hist} or \texttt{mod}. A scan's input is wrapped in \texttt{opaque},
  because Futhark 0.28.1's SOAC fusion fails on a scan of a scan with an
  internal compiler error (measured on every pair of the three operators).
\end{itemize}
All element operators are total. \texttt{i32} arithmetic wraps, and there is
no division, which traps on zero.

The e-graph's rules are exact identities of wrapping \texttt{i32} arithmetic.
They are commutativity of \texttt{+}, \texttt{*}, \texttt{min} and
\texttt{max}; associativity of \texttt{+}, \texttt{min} and \texttt{max};
$a-b=a+(-b)$; and $-(-a)=a$. For arrays they are commutativity of
\texttt{map2}, and \texttt{map2 min a a = a} (and the same for \texttt{max}).
Each is generated in every context.

\paragraph{Worked example: histogram.} The published program is
\texttt{hist (+) 0 k (map (\%k) is) vs}. As a gene of the kingdom
(\src{src/futhark_sig.rs:200}) it is
\begin{lstlisting}
(hist.add.n (tab (mod.k (get.n is j1))) vs)
\end{lstlisting}
which renders, by the rows' spellings, as
\texttt{hist (+) 0i32 k (map i64.i32 (tabulate n (\textbackslash j1 -> is[j1] \% i32.i64 k))) vs}.
Here \texttt{j1} is a terminal only because the \texttt{tabulate} binds it, and
\texttt{get.n} takes an \ty{ix.n}, so the read is in range.

\paragraph{Worked example: pathfinder.} The published kernel uses \texttt{map5},
\texttt{rotate} and edge branches. The kingdom writes the same algorithm with
clamped neighbours (\src{src/futhark_sig.rs:145}):
\begin{lstlisting}
(loop1.rows (row.rows.cols walls zero)
  (tab (add (get.cols (row.rows.cols walls i1) j2)
            (min (min (get.cols s1 j2) (get.cols s1 (prev j2)))
                 (get.cols s1 (next j2))))))
\end{lstlisting}
At an edge, the clamped neighbour is the cell itself, which the minimum already
contains, so the values are the same. This form turns out to be faster than the
published port (Section~\ref{sec:study}).

\subsection{The pixel family}
A picture is a \texttt{map} over pixels of a scalar program. The program is a
chain of \texttt{let}s and bounded \texttt{loop}s over \texttt{f32}, ending in
\texttt{[r, g, b]}. Scope is a type (\ty{f32.s}$k$ has $v_0\ldots v_{k-1}$ in
scope). A binding moves the scope. Trip counts are literals of at most 64, and
loops nest at most two deep (\src{kingdoms/futhark_shapes/gen_tables.py:1130-1222}).
An instance is declared by its WGSL source: the uniform blocks become inputs,
and the source's float literals join the constant pool. The hand-written
references draw the WGSL originals' pictures up to fast-math differences
(volumepassage: 13 of 16,384 pixels differ by more than 1/255). Evolution from
random did not draw any target picture, and the study is parked
(\src{docs/reports/2026-10-27_image_study_parked.md}). We report nothing more
from this family.

\subsection{The deprecated 1-D kingdom}
The first Futhark kingdom (\src{kingdoms/futhark_true/}, 4 types and 38 rows,
\src{src/futhark_kingdom.rs}) made \emph{shape a value, not a type}: \ty{arr.f32}
was one type whatever its length. That kept the kingdom small, but it could not
say that two arrays' shapes are related, so a \texttt{matmul} was not
expressible. Two findings from it remain useful. Its e-graph diverged when
\texttt{map id} was combined with a fusion rule (5, 7, 13, 55 and then 1,357
e-nodes) until the unit laws of composition were added. And one legal program
made Futhark 0.28.1's SOAC fusion loop forever; capping that loop
(\texttt{--fuse=10}) compiles it in 0.12\,s. The kingdom is superseded and is not
used in any experiment here.

\subsection{Reading real programs: the reader and the ring canvas}
\label{sec:reader}
\src{src/futhark_reader.rs} reads Futhark source back into a tree of a
signature kingdom. It parses a subset of the language, inlines \texttt{let}s and
helpers, and matches each expression against the rows' own spellings. A
construct it cannot express is recorded as a \emph{gap}, never skipped. As a
round trip, 1,600 rendered random genes (400 per instance) read back to the
same program text.

When a gene is written, its unread cells are filled with random symbols. A
mutation that grows the tree then exposes those cells, so random symbols come
back within a few generations. Andrew's \emph{ring canvas} fills the unread cells
with real code instead (\src{src/futhark_rings.rs:183}). Read programs are
converted to Karva and their coding symbols are appended, by type, to rings of
head symbols, terminals and constants. A new gene is a window onto each type's
rings. The pool holds the benchmarks' own functions and programs that Sonnet
wrote on 60 unrelated topics (it was never told the task). For histogram 48 of 60
of those programs could be read, and for pathfinder 45. The published program's
own file is left out. Measured on 1,400 fresh histogram genes, the canvas cut
the share of programs whose output is all zeros from 40.8\,\% to 13.8\,\%, and
raised the number of distinct outputs from 354 to 530
(\src{docs/PLAN_futhark_fitness_fixes.md}, section H).

%======================================================================
\section{Fitness evaluation}
\label{sec:fitness}

\subsection{Cases}
Each benchmark's cases are random inputs paired with the published program's
exact outputs. They are split into training, held-out and test cases. Training
and held-out cases are fitness inputs (the held-out cases form the quality
group's validation column). Test cases are run but are never fitness inputs.
After the fixes (Section~\ref{sec:faults}), histogram has 11 training, 4
held-out and 4 test cases. Pathfinder has 13, 6 and 4. The training sets include
\emph{stepping stones}: cases that a partly correct program can solve exactly,
such as histogram indices that are already in range, or $k=1,2,3$
(\src{examples/futhark_problem_cases.rs}).

\subsection{Quality inputs and the staged HFF}
Each case contributes its own inputs on $[0,1]$, where 0 is exact
(\src{src/futhark_sig.rs:1183}, \texttt{QUALITY\_PARTS};
\src{:1244}, \texttt{quality\_inputs}):
\begin{itemize}
\item the share of cells that are not exact;
\item $(1-r)/2$ for Spearman's and for Pearson's correlation with the correct
  output (left out when the correct output is constant);
\item the share of cells that are off by more than 50\,\% of the correct value's
  magnitude, and the share off by more than 1\,\%.
\end{itemize}
This formula was chosen by the slope check (Section~\ref{sec:slope}), from every
combination of 13 measured parts. Nothing is averaged before HFF: each input
keeps its own place in TrueNorth's angle,
$\theta=\arccos\!\bigl(1-\tfrac1m\sum x_i^2\bigr)$, which is computed by the hff
library~\cite{morgan2026hff}.

The fitness is staged, by Andrew's rule: quality first, then performance
(\src{src/runner.rs:97}, \texttt{staged\_hff}). In stage 1 the fitness is
TrueNorth over the quality inputs, and no program is timed. When the best
individual is exact on every training and held-out case, the whole population
moves to stage 2, and it never moves back. Stage 2 is the two-level HFF:
TrueNorth over $(\theta_q/\tfrac\pi2,\ \theta_s/\tfrac\pi2)$, where $\theta_s$ is
the angle of the speed inputs. Stage 2 is \emph{gated by correctness}. An exact
program scores the two-level angle, in $[0,\pi/2)$. An inexact one scores
$\pi/2+\theta_q$. So no inexact program, however fast, ranks above an exact one.
We added this gate after a run in which wrong but fast programs outranked an
exact seed (Section~\ref{sec:reruns}).

\subsection{Speed inputs and the verdict}
A program's speed inputs are its \emph{timing slots}
(\src{src/runner.rs:123}). There are 11 timed pairs against the published
program, kept by the program's text, so a program scored again in a later
generation reads the same pairs. Each pair is its own input, $r/(1+r)$, where
$r$ is the candidate's GPU time over the published program's. In a pair, the
two programs alternate run by run on one upload of the timing input, at least
1\,s a side, and each side takes its fastest run
(\src{src/runner.rs:220}; \src{src/futhark_batch.rs:98}). Before any timed run
an A/A gate times the published program against itself, and every A/A ratio
must lie within $\pm5\,\%$. A champion counts as \emph{faster} only if the median
of its 11-pair re-time is below the least A/A ratio \emph{and} at least 9 of its
11 pairs are below 1. That is a sign test with one-sided $p=0.033$
(\src{src/runner.rs:202}). Before scoring, the host computes each program's
exact node executions on the case sizes (\src{src/futhark_sig.rs:877}). A
program above the device's per-submit cap stops the run as misconfigured. A
program above the timing cap is not timed.

\subsection{The slope check}
\label{sec:slope}
A run may start only after a slope check of the same problem, program size,
cases and quality formula has passed (\src{src/runner.rs:388};
\src{examples/futhark_evolve.rs:1372}). The check scores the reference gene,
every one-edit neighbour of it, and 1,400 random genes on the GPU. It passes
only if (i) the median inexact neighbour is better than at least 95\,\% of the
random genes, and (ii) Spearman's correlation between fitness and tree edit
distance to the reference~\cite{zhang1989ted} is positive. Lower fitness is
better, so a positive correlation means that programs nearer the reference
score better. The 95\,\% bar was 99\,\% when the check was designed. It was
lowered after the first three measurements, with Andrew's approval, because
ties at outputs that random genes also produce cap the reachable figure at about
99.4\,\% on histogram (\src{docs/PLAN_futhark_fitness_fixes.md}, section A).

\subsection{The matrix family's fitness}
Correctness is judged against the \emph{reference kernel's} value on a suite of
independent input sets. The comparison is bit for bit (signed zeros told
apart) on element-wise paths, and exact (no tolerance) where there is a
reduction (\src{src/futhark_shapes_interp.rs:281-297}). The fitness that the code
computes today is the two-level HFF. Its quality inputs are the disagreement on
each input set, on HFF's log scale. Its speed inputs are traffic, allocations
and depth, each as $x/(x+x_{\mathrm{ref}})$
(\src{src/futhark_shapes_evolve.rs:1091-1106}).

%======================================================================
\section{GPU implementation}
\label{sec:gpu}

\subsection{Three evaluation paths}
The signature runner (\src{examples/futhark_evolve.rs}) can evaluate a
generation in three ways. \emph{Single} compiles each candidate with the
WebGPU compiler, links it against our wgpu shim and runs it in its own process.
\emph{Batch} makes each new program of a generation one entry point of a single
Futhark library, run by one driver process. If a program fails to compile, the
batch is split in halves until that program stands alone; it is logged and
scored worst (\src{src/futhark_batch.rs:1-33}). \emph{Rust} compiles the same
batch with our fork's \texttt{rust-wgpu} backend into one Rust module. That
module holds the WGSL kernels and a host over wgpu-rs. \texttt{rustc} builds it
into a dylib, which the runtime \texttt{futhark-rust-rt} loads and runs inside
this process on one device: no C, no shim and no process per candidate
(\src{src/futhark_rust.rs:1-12}, \src{:128}, \src{:257}). The kernels are
launched on the device (\src{futhark-rust-rt/src/program.rs:252}), and each pass
is timed with GPU timestamps (\src{futhark-rust-rt/src/ctx.rs:114-119}).

\subsection{The compiler fork}
The fork is a series of nine patches on Futhark's \texttt{webgpu} branch at
\texttt{14e9ccb} (\src{futhark-rust-backend/README.md}). Six are needed for the
backend and for the matrix kernel (the backend itself; SetMem of register memory;
names; register arrays sized by a kernel constant; a boolean kernel constant
that was declared \texttt{i64}, an upstream defect). Three came from the
study's compile failures (Table~\ref{tab:fixes}). The runtime reimplements the
C runtime's memory and free list, pipeline cache, launches, tuning and copies.
It differs from the C runtime in no way that changes a value
(\src{docs/PLAN_futhark_rust_wgpu.md}, section 3).

\begin{table}[t]
\caption{The rust path against \texttt{futhark c} and the shim, Apple M3 Max,
wgpu 30.0.1 on Metal (\src{docs/PLAN_futhark_rust_wgpu.md}, section 7). GPU
$\mu$s is the fastest of 20 runs; k-means sums f32 with atomics and does not
repeat from run to run on either GPU path.}
\label{tab:rust}
\small
\begin{tabular}{@{}lrrrr@{}}
\toprule
benchmark & cases & $=$\,C (bits) & rust $\mu$s & shim $\mu$s\\
\midrule
reduce     & 5 & 5 & 8.5  & 27.6\\
scan       & 5 & 5 & 16.8 & 17.4\\
histogram  & 5 & 5 & 23.3 & 11.0\\
pathfinder & 8 & 8 & 56.2 & 57.8\\
k-means    & 5 & 1 & 43.5 & 43.7\\
\bottomrule
\end{tabular}
\end{table}

On the four integer benchmarks the rust path is bit-identical to
\texttt{futhark c} on every case (Table~\ref{tab:rust}). Over three generations
of evolution its ledger (every fitness, every text, every quality input)
matched the batch path's with 0 mismatches on all four problems, and it scored
0--5\,\% more evaluations per second.

\subsection{Where each stage runs}
Table~\ref{tab:where} gives, with file and line, where each stage runs in the
two runners. fuller must not depend on phylu, so the signature runner has its
own host genetic algorithm. In phylu's \texttt{shape\_typed} runner, the
population lives in the typed core on the device, and the host scores it.

\begin{table*}[t]
\caption{Where each stage runs. \D{} device, \Hh{} host. ``Signature'' is fuller's
\texttt{examples/futhark\_evolve.rs} (the benchmark study); ``Matrix (GPU population)''
is phylu's \texttt{src/evolve/shape\_typed.rs} with fuller's scorer. Paths are
relative to each repository.}
\label{tab:where}
\small
\begin{tabularx}{\textwidth}{@{}lcYcY@{}}
\toprule
stage & \multicolumn{2}{c}{Signature (fuller)} & \multicolumn{2}{c}{Matrix, GPU population (phylu + fuller)}\\
\midrule
random genes & \Hh & \src{futhark_evolve.rs:790} $\to$ \src{src/futhark_sig.rs:1007} & \D & typed core init (\src{phylu src/evolve/shape_typed.rs:10-17})\\
expression & \Hh & \src{futhark_evolve.rs:828} & \Hh{}\,(+\D{} walk check) & \texttt{score\_lenient} \src{futhark_shapes_evolve.rs:1012}; walk check \src{shape_typed.rs:566}\\
mutation, crossover, marriage & \Hh & \src{futhark_evolve.rs:838}, \src{:867}, \src{:882} & \D & \texttt{dev.vary} \src{shape_typed.rs:746}\\
e-graph variants & \Hh & egglog, \src{src/futhark_sig.rs:1444} & \Hh & \texttt{saturate\_tree} \src{shape_typed.rs:630}\\
program building & \Hh & render \src{futhark_sig.rs:423}; fork + \texttt{rustc} \src{futhark_rust.rs:128} & \Hh & interpreter; compile only for measured elites \src{futhark_shapes_gpu.rs:110}\\
program execution & \D & \src{futhark-rust-rt/src/program.rs:252} & \Hh & interpreter \src{futhark_shapes_interp.rs:331}\\
timing & \D & GPU timestamps \src{futhark-rust-rt/src/ctx.rs:114}; pairs \src{src/runner.rs:220} & -- & static cost \src{futhark_shapes_cost.rs:263} (no timing)\\
fitness evaluation & \Hh & \src{futhark_sig.rs:1244}, \src{src/runner.rs:97} & \Hh & \src{futhark_shapes_evolve.rs:1104}; written to \D{} \src{shape_typed.rs:488}\\
selection & \Hh & tournament of 3, \src{futhark_evolve.rs:2362} & \D & device selection kernel on the host's fitness\\
\bottomrule
\end{tabularx}
\end{table*}

The matrix family's evaluation stays on the host because its contract is
bit-exactness. The GPU's backend contracts multiply-add into an FMA and the
interpreter does not: on the A7 product, 758 of 4,096 elements agreed bit for
bit, and the rest differed by up to $3.7\times10^{-4}$
(\src{docs/measurements/2026-10-27_futhark_shapes_A7.md}).

%======================================================================
\section{Experiments}
All runs were on one machine, an Apple M3 Max, under Metal. The A7 run used
Futhark 0.28.1's OpenCL backend. The benchmark runs used our wgpu path (wgpu
30.0.1). Expectations were written before the runs; where a rule was changed
after a result was seen, we say so.

\subsection{The matrix kernel A7}
\label{sec:a7}
\paragraph{Setup.} A7 computes the product of two matrices twice, once directly
and once the transposed way, and adds them:
\texttt{map2 (add) (A.B) ((B'.A')')}. The sizes were $n=m=p=64$; the population
was 32, run for 20 generations, starting from random genes plus the reference
(\src{docs/measurements/2026-10-27_futhark_shapes_A7_learning.md}).

\paragraph{Result.} The best program is \texttt{map2 (add) (A.B) (A.B)}. The
e-graph's transposed-product rule rewrote the second operand into the first, and
Futhark's common subexpression elimination removed the duplicate
(Table~\ref{tab:a7}). On Apple's OpenCL backend, the reference does not run: its
kernel is too large to build (\texttt{newCompute\-Pipeline\-State failed}, error
$-48$). The winner runs. Through our rust path both programs compile and run,
and both agree with \texttt{futhark c} on 1,157 of 4,096 values, within 6 ulp
(\src{docs/PLAN_futhark_rust_wgpu.md}, section 7). So ``the reference cannot
run'' describes one backend, not the program.

\begin{table}[t]
\caption{A7 at $n=m=p=64$, Apple OpenCL over Metal
(\src{docs/measurements/2026-10-27_futhark_shapes_A7.md}).}
\label{tab:a7}
\small
\begin{tabular}{@{}lrrrrrl@{}}
\toprule
program & traffic & dev.\ allocs & bytes & kernels & tiled & runs\\
\midrule
reference & 45,056 & 15 & 188,416 & 10 & yes & no\\
winner    & 20,480 & 9  & 106,496 & 5  & no  & yes\\
\bottomrule
\end{tabular}
\end{table}

\begin{figure}[t]
\centering
\begin{tikzpicture}
\begin{axis}[width=\columnwidth,height=4.2cm,xlabel={generation},ylabel={correct of 32},
  ymin=0,ymax=34,xmin=0,xmax=20,legend style={font=\scriptsize,at={(0.98,0.5)},anchor=east},
  tick label style={font=\scriptsize},label style={font=\scriptsize}]
\addplot[evo,thick,mark=*,mark size=1.2pt] table[x=generation,y=correct]{\dat{a7_host.dat}};
\addlegendentry{host population, $64^3$}
\addplot[rnd,thick,mark=square*,mark size=1.2pt] table[x=generation,y=correct]{\dat{a7_gpu.dat}};
\addlegendentry{device population, $8^3$}
\end{axis}
\end{tikzpicture}
\caption{A7: individuals bit-identical to the reference over the input suite.
Host run: \src{docs/measurements/2026-10-27_futhark_shapes_A7_learning.md}. Device
run (phylu \texttt{shape\_fit}):
\src{docs/measurements/2026-10-27_futhark_shapes_A7_gpu_population.md}.}
\label{fig:a7}
\end{figure}

\paragraph{Learning.} In the host run, the share of correct individuals rose
from 1 of 32 at generation 0 to 32 of 32 at generation 4. After that it moved
between 28 and 32 (Figure~\ref{fig:a7}). The population converged on the
untiled, runnable form: from generation 2 on, the device refused no measured
candidate and the compiler tiled none. The cost moved only once. The e-graph improved one elite
at generation 0, and no generation after it improved any. The e-graph's rules
are identities, so they can move only the cost axis. GEP's operators change the
program, so they explore correctness. With the population on the device
($n=m=p=8$), the correct count moved between 5 and 12 of 32 and never settled,
although the best individual was correct. The fitness recorded in that run is the
earlier two-band form: its logged best, 0.15313, equals
$0.49\cdot320/(320+704)$. The code now computes the two-level HFF instead
(\src{src/futhark_shapes_evolve.rs:1104}), so that run predates the current
fitness.

\subsection{The benchmark study}
\label{sec:study}
\paragraph{Question and hypotheses.} Given only examples of a published
benchmark's inputs and outputs, does evolution produce the program exactly
(H1)? Does it do so sooner than random search at the same budget (H2)? Starting
from the published program, does it find a faster equivalent (H3)? The
hypotheses and expectations are written in \src{docs/STUDY_evolving_benchmarks.md}.
We expected H1 to hold on reduce and scan and to be likely on histogram, and H3 to
be unlikely on reduce and scan and possible on histogram and pathfinder.

\paragraph{Setup.} Four benchmarks: reduce, scan, histogram (Futhark's
\path{micro/reduce_by_index}) and Rodinia pathfinder. Three regimes: R,
evolution from random, 3 seeds; C, random search (a new random gene for each
evaluation), 3 seeds; and P, starting from the published program and searching
for a faster one, 1 seed for 60 minutes. The budget was 1,400 evaluations, with
population 24, 2 elites and tournaments of 3. The operators were point mutation,
homologous crossover, marriage (typed subtree crossover between parents whose
trees differ), the e-graph pump, and the saturated fold. Each P run was preceded
by an A/A gate and a control in which a slowdown was planted on purpose.

\begin{table}[t]
\caption{The benchmark study, before the fixes
(\src{docs/reports/2026-10-27_STUDY_evolving_benchmarks_summary.md}). ``First exact'' gives
the evaluation at which a seed was first exact; ``best'' is the best-ever quality
angle (0 exact) of the three seeds.}
\label{tab:study}
\small
\begin{tabular}{@{}lll@{}}
\toprule
benchmark & R (evolution) & C (random search)\\
\midrule
reduce (2 nodes)     & exact 3/3: 19, 118, 119  & exact 3/3: 19, 127, 60\\
scan (2 nodes)       & exact 3/3: 57, 108, 41   & exact 3/3: 42, 46, 28\\
histogram (7 nodes)  & 0/3; best 1.124--1.147    & 0/3; best 1.120--1.137\\
pathfinder (24 nodes)& 0/3; best 1.015--1.207    & 0/3; best 1.076--1.105\\
\midrule
\multicolumn{3}{@{}l}{H3 (P): reduce undecided (A/A gate 21/22); scan matched 1.0000;}\\
\multicolumn{3}{@{}l}{histogram matched 1.0046; pathfinder 0.9705, 11 of 11 pairs faster.}\\
\bottomrule
\end{tabular}
\end{table}

\paragraph{Results.} Evolution was no better than random search
(Table~\ref{tab:study}). H1 held only for the two-node programs, and random
search found those as fast or faster: a random gene of program size 8 expresses
\texttt{(sum.n xs)} about once in 20 to 130 draws. Evolution's own exact reduce
programs were bloated equivalents (14 and 10 nodes) that ran 4.6 and 2.4 times
slower. On histogram, random search got as close as evolution or closer. On
pathfinder, two evolution seeds got closer than any random search seed, and the
third did not move from 1.207.

\paragraph{The pathfinder ``beat''.} P's champion was 0.9705 of the published
program's time, with 11 of 11 pairs faster. But its starting point, the reference
gene (the published algorithm as the kingdom writes it, never evolved), re-timed
at 0.9473 and 0.9434, also 11 of 11. Evolution returned an e-graph variant that
was about 2.5\,\% \emph{slower} than where it started. Selection had preferred it
because of noise in the 11 in-run slots (0.940 against 0.955, where the re-times
give 0.970 against 0.945). The 5\,\% gain is real and was repeated in 22 of 22
pairs. It comes from the kingdom's spelling, which clamps indices instead of
branching and uses \texttt{tabulate} with direct indexing instead of
\texttt{map5} over two \texttt{rotate}s. It does not come from search.

\subsection{The faults}
\label{sec:faults}
During the study we found ten faults in our own method, fixed each one, and
reran or voided every result it affected (the summary's ``methods findings''
table). Among them: timing a 10--30\,$\mu$s kernel measured the GPU's clock state
(A/A ratios in three clusters, 0.30, 1.00 and 3.3), so micro benchmarks now have
their own timing input; the first rule for ``faster'' gave a false win to the
published program itself (median 0.985, but only 8 of 11 pairs faster), so the
sign test was added; random search's champion had been reported from its last
batch rather than from the hall of fame; and an input of $n=2^{24}$ exceeded
wgpu's dispatch limit of 65,535 workgroups.

Andrew's premise was that random search beats evolution only when the fitness is
broken. An independent review then read the runner, the scorer and TrueNorth, and
re-implemented the quality inputs. The re-implementation reproduced the
ledgers' fitness exactly (\src{docs/reviews/2026-10-28_futhark_evolve_faults.md}).
It found 22 faults: 6 break evolution, 11 distort measurement, and 5 are minor
(Table~\ref{tab:faults}). The central one is fault 1. Over all the distinct
programs the study evaluated, fitness got \emph{better} as tree distance from
the reference grew: Spearman $-0.68$ on histogram (6,588 programs) and $-0.52$
on pathfinder (6,027). The target's one-edit neighbours scored in the same band
as unrelated programs. On pathfinder, dropping one neighbour from the
\texttt{min} gives a program right in structure and close in value, which
scored 1.045. That is worse than a 46-node program wrong in 87\,\% of cells
(1.015). Selection therefore moved the population away from the target.

\begin{table*}[t]
\caption{The 22 faults of the review
(\src{docs/reviews/2026-10-28_futhark_evolve_faults.md}) and the fix built for each
(\src{docs/PLAN_futhark_fitness_fixes.md}, sections A--G). B = breaks evolution, M =
distorts measurement, m = minor. Three were suspected rather than confirmed (17, 19, 22).}
\label{tab:faults}
\footnotesize
\begin{tabularx}{\textwidth}{@{}rlYY@{}}
\toprule
\# & & fault & fix (built unless stated)\\
\midrule
1 & B & fitness does not reward closeness: Spearman(fitness, distance) $-0.68$, $-0.52$ & the slope check, required before any run (A)\\
2 & B & four of seven inputs are step functions, three near-copies & new quality formula chosen by the check (B)\\
3 & B & error scaled by the output's range rewards matching the level & shape inputs (Spearman, Pearson); magnitude, not range (B)\\
4 & B & histogram's whole-range values wrap; no case gives partial credit & stepping-stone cases (C)\\
5 & B & degenerate cases (one row, one bin) reward trivial programs & more normal cases (C)\\
6 & B & behaviour diversity collapses (23 texts, about 5 outputs) while text diversity is reported & elites, partners and pool differ in output; distinct outputs reported (D)\\
7 & M & elites re-scored and counted (8.3\,\% of R's budget) & an evaluation is a program not scored before (E)\\
8 & M & repeats count as evaluations & as 7 (E)\\
9 & M & pump and fold cannot change stage-1 fitness, and repeat themselves & pump only when the best changed; repeats not counted (F)\\
10 & M & R and C share generation 0 & C has its own seed stream (E)\\
11 & M & held-out cases are in the fitness & a test split that selection never sees (G)\\
12 & M & \texttt{--rescore} uses the old verdict rule & one verdict function (G)\\
13 & M & batched timing fails the A/A gate; selection on in-run slots chose noise; the floor is an extreme value; pathfinder timed on a quality case & interleaved pairs (Section~\ref{sec:fixes}); the floor and pathfinder's timing input are \emph{not} changed\\
14 & M & legal programs the compiler cannot build are scored worst & patches 0007--0009 (Section~\ref{sec:fixes})\\
15 & M & champion fields from the last batch & from the hall of fame's rank 1 (G)\\
16 & M & hall of fame mixes stage-1 and stage-2 fitness & one hall of fame per stage (G)\\
17 & M & seed 1 on a different binary and path & one binary per comparison (G)\\
18 & m & bloat under neutral drift (8--17 to 35--124 nodes) & symptom of 1, 2, 6; no length penalty\\
19 & m & marriage favours swaps near the root & uniform over fitting pairs (F)\\
20 & m & crossover drops its second child; always followed by mutation & both kept; mutation a recorded choice (F)\\
21 & m & partners need only a different text & differ in output (D)\\
22 & m & 24 individuals, 58 generations: a small search & not changed\\
\bottomrule
\end{tabularx}
\end{table*}

\paragraph{The slope check, measured.} Table~\ref{tab:slope} shows the check on
the study's fitness and on the chosen one. Under the study's formula, histogram
failed on the old cases (A1) and passed only once the stepping stones were added
(A3). Pathfinder failed with every set of cases (A2, A4, A5). The chosen formula
passes both benchmarks (B1). It also no longer rewards the all-zeros output:
68.6\,\% (histogram) and 84.1\,\% (pathfinder) of non-zero wrong random outputs
now score better than it, against 28.8\,\% on histogram under the study's
formula. Graded per-cell inputs were measured and rejected: on histogram, cells
holding unrelated large numbers drowned the exactness signal, and on pathfinder
a constant at the right level beat the near-neighbours
(\src{docs/measurements/2026-10-29_futhark_slope_check.tsv}, 82 formula and case
combinations).

\begin{table}[t]
\caption{The slope check on the GPU, 1,400 random genes
(\src{docs/PLAN_futhark_fitness_fixes.md}, ``Measurements of the check'';
\src{results/2026-10-27_31/fcheck/}). ``Beats'' = share of random genes the median
neighbour beats; pass needs $\ge95\,\%$ and $\rho>0$.}
\label{tab:slope}
\small
\begin{tabular}{@{}lllrrl@{}}
\toprule
 & benchmark & formula, cases & beats & $\rho$ & verdict\\
\midrule
A1 & histogram  & study, old  & 98.86\,\% & $-0.080$ & fail\\
A2 & pathfinder & study, old  & 92.71\,\% & $+0.437$ & fail\\
A3 & histogram  & study, new  & 98.50\,\% & $+0.370$ & pass\\
A4 & pathfinder & study, new  & 93.00\,\% & $+0.421$ & fail\\
A5 & pathfinder & study, new+4 & 92.93\,\% & $+0.420$ & fail\\
B1 & histogram  & chosen, new & 97.36\,\% & $+0.079$ & pass\\
B1 & pathfinder & chosen, new & 99.00\,\% & $+0.403$ & pass\\
\bottomrule
\end{tabular}
\end{table}

\subsection{Reruns on the corrected fitness}
\label{sec:reruns}
\paragraph{Setup.} Fuller \texttt{72c37606}, fitness of section B, both slope
checks passed. Generation 0 was ring-canvas genes. Population 24, budget 1,400
new programs, a 30-minute bound per run, \texttt{--eval rust}, point mutation, and
seeds 1 and 2. C drew from its own seed stream. We expected R's best to be
better than C's in at least 2 of 3 seeds per benchmark, and we did not expect an
exact program (\src{docs/PLAN_futhark_fitness_fixes.md}, ``Expectations'').
Two seeds were run, not three.

\begin{table}[t]
\caption{Best HFF (0 exact, lower better) at the same number of programs, the
shorter run's (\src{docs/reports/2026-10-29_futhark_reruns.md};
\src{results/2026-10-27_31/frerun2/point/}). The histogram R runs reached the
30-minute bound before 1,400 programs.}
\label{tab:reruns}
\small
\begin{tabular}{@{}lrrrrr@{}}
\toprule
benchmark, seed & programs & R best & C best & R median & C median\\
\midrule
histogram, 1  & 689   & \textbf{1.0584} & 1.1346 & 1.107 & 1.293\\
histogram, 2  & 1,218 & \textbf{0.9513} & 1.0608 & 0.953 & 1.286\\
pathfinder, 1 & 1,400 & \textbf{0.9834} & 1.0000 & 0.984 & 1.243\\
pathfinder, 2 & 1,400 & \textbf{0.9536} & 1.0000 & 0.964 & 1.261\\
\bottomrule
\end{tabular}
\end{table}

\begin{figure*}[t]
\centering
\pgfplotsset{every axis/.style={width=0.27\textwidth,height=4.0cm,
  tick label style={font=\scriptsize},label style={font=\scriptsize},
  title style={font=\scriptsize},ymin=0.9,ymax=1.32,xmin=0,xmax=1420}}
\begin{tikzpicture}
\begin{axis}[title={histogram, seed 1},xlabel={programs scored},ylabel={HFF}]
\addplot[evo,thick] table[x=evaluations,y=best]{\dat{rerun_histogram_R_s1.dat}};
\addplot[evo,dashed] table[x=evaluations,y=median]{\dat{rerun_histogram_R_s1.dat}};
\addplot[rnd,thick] table[x=evaluations,y=best]{\dat{rerun_histogram_C_s1.dat}};
\addplot[rnd,dashed] table[x=evaluations,y=median]{\dat{rerun_histogram_C_s1.dat}};
\end{axis}
\end{tikzpicture}%
\begin{tikzpicture}
\begin{axis}[title={histogram, seed 2},xlabel={programs scored},yticklabels={}]
\addplot[evo,thick] table[x=evaluations,y=best]{\dat{rerun_histogram_R_s2.dat}};
\addplot[evo,dashed] table[x=evaluations,y=median]{\dat{rerun_histogram_R_s2.dat}};
\addplot[rnd,thick] table[x=evaluations,y=best]{\dat{rerun_histogram_C_s2.dat}};
\addplot[rnd,dashed] table[x=evaluations,y=median]{\dat{rerun_histogram_C_s2.dat}};
\end{axis}
\end{tikzpicture}%
\begin{tikzpicture}
\begin{axis}[title={pathfinder, seed 1},xlabel={programs scored},yticklabels={}]
\addplot[evo,thick] table[x=evaluations,y=best]{\dat{rerun_pathfinder_R_s1.dat}};
\addplot[evo,dashed] table[x=evaluations,y=median]{\dat{rerun_pathfinder_R_s1.dat}};
\addplot[rnd,thick] table[x=evaluations,y=best]{\dat{rerun_pathfinder_C_s1.dat}};
\addplot[rnd,dashed] table[x=evaluations,y=median]{\dat{rerun_pathfinder_C_s1.dat}};
\end{axis}
\end{tikzpicture}%
\begin{tikzpicture}
\begin{axis}[title={pathfinder, seed 2},xlabel={programs scored},yticklabels={}]
\addplot[evo,thick] table[x=evaluations,y=best]{\dat{rerun_pathfinder_R_s2.dat}};
\addplot[evo,dashed] table[x=evaluations,y=median]{\dat{rerun_pathfinder_R_s2.dat}};
\addplot[rnd,thick] table[x=evaluations,y=best]{\dat{rerun_pathfinder_C_s2.dat}};
\addplot[rnd,dashed] table[x=evaluations,y=median]{\dat{rerun_pathfinder_C_s2.dat}};
\end{axis}
\end{tikzpicture}
\caption{HFF against programs scored on the corrected fitness: best so far (solid)
and the generation's median (dashed), evolution (blue) against random search (red).
Data: \src{papers/futhark_multityped_gep_data/rerun_*.dat}, taken from each run's
\texttt{per\_generation.tsv}.}
\label{fig:reruns}
\end{figure*}

\paragraph{Results.} Evolution reached a lower best HFF than random search in all
4 pairs, at equal numbers of programs scored (Table~\ref{tab:reruns},
Figure~\ref{fig:reruns}). Evolution's median fell during every run. Random
search's median stayed between about 1.24 and 1.29. On histogram, evolution's best
after a run cut short (1.058 at 689 programs; 0.951 at 1,218) was better than
random search's best after all 1,400 (1.093 and 1.061). Random search's best on
pathfinder stopped at 1.0000 in both seeds. No run found an exact program, and
none reached stage 2. The histogram R runs reached the 30-minute bound because
of compile failures (74 to 133 a run), which made batches split in halves. The
compiler patches that remove those failures were not yet in these runs.

\paragraph{Ring walk.} Ring walk replaces point mutation by moving a type's
reading frame along its ring. It did worse. On histogram, seed 1 ran 16,149
generations and produced only 858 new programs among 386,742 repeats, and its
best never left the generation-0 value (1.1346). On pathfinder it learned
(1.140 to 1.027), but less than point mutation (0.983). Point mutation stays.

\begin{table}[t]
\caption{Operator rates. Top: children in the reruns, 2 seeds
(\src{docs/PLAN_futhark_fitness_fixes.md}, section I). Bottom: 700 children per
operator, one application each, ring canvas
(\src{results/2026-10-27_31/fmeasure/fix/*/measure_summary.tsv}).}
\label{tab:ops}
\small
\begin{tabular}{@{}lrrr@{}}
\toprule
operator & new output & better than parent & same output\\
\midrule
point mutation       & 39--60\,\% & 4.4--7.9\,\% & 12--29\,\%\\
crossover+mutation   & 44--59\,\% & 1.5--3.2\,\% & 8--24\,\%\\
marriage             & 29--48\,\% & 3.1--6.3\,\% & 34--58\,\%\\
crossover alone      & 15--22\,\% & 1.7--5.7\,\% & 70--79\,\%\\
pump (e-graph)       & 0\,\%      & 0\,\%        & 100\,\%\\
saturated fold       & 0\,\%      & 0\,\%        & 99--100\,\%\\
ring walk            & 0.1--1.7\,\% & 1.5--2.7\,\% & 53--56\,\%\\
\midrule
point, hist / path   & 45.9 / 44.9\,\% & 4.0 / 15.3\,\% & 42.6 / 43.0\,\%\\
inversion            & 45.9 / 56.1\,\% & 3.7 / 18.1\,\% & 32.3 / 24.7\,\%\\
transposition        & 46.1 / 55.1\,\% & 3.7 / 14.4\,\% & 34.6 / 28.1\,\%\\
\bottomrule
\end{tabular}
\end{table}

\paragraph{Operators.} Table~\ref{tab:ops} gives the measured rates. The pump and
the fold never change stage-1 fitness. That is by design: their rules are
identities, so they matter only once programs are timed. Inversion and
transposition (IS and RIS within one type's head) were measured alone, 700
children each, but have not yet been tested inside runs. On pathfinder they
give more new outputs than point mutation, and inversion more improvements;
on histogram all three are similar.

\paragraph{Programs from a language model.} In one generation-0 set, a program
written by Fable (a language model) was one of the individuals. All 8 such runs
stopped at evaluation 1, because Fable's program was already exact. On
histogram it was 70 to 72 times slower than the published program: it
tabulates over the whole input once for each bin. On pathfinder it was 0.94 to
0.96 of the published time, the same algorithm written with \texttt{map2} over
whole rows. These are a model's programs, not evolution's.

\paragraph{Improvement run, and the correctness gate.} An improvement run from
Fable's exact histogram program showed a fault: under the two-level HFF, wrong
but fast mutants outranked the exact, slow seed. One mutant had quality angle
1.29 and speed ratio 0.002 and scored 0.848. Over 7 generations the best was
inexact in every generation. The run was stopped and the stage-2 gate
described in Section~\ref{sec:fitness} was built. With the gate in place, the
best stayed at the exact seed for 18 generations and no faster exact program
appeared (\src{results/2026-10-27_31/frerun2/improve/histogram_R_s1_gate/}). A
further run started from Fable's mutants with Fable's program itself left
out. Its best HFF fell from 1.2644 to 1.1024 ($-13\,\%$) over 23 generations
(445 programs), and its median from 1.2926 to 1.1358. It found no exact
program.

\subsection{Compiler and timing fixes}
\label{sec:fixes}
At the end of the reruns four bugs were open
(\src{docs/reports/2026-10-29_futhark_reruns.md}). All four are fixed and tested
(Table~\ref{tab:fixes}). The three paths now agree. On five programs (the
published histogram, its planted slowdown, Fable's program, a per-bin sum and a
constant) the quality outputs and inputs are bit-identical on all three paths,
and the median speed ratios agree; Fable's program, for example, is 70.75, 70.64
and 70.47. Generation 0 is byte-identical in the three ledgers.

\begin{table}[t]
\caption{Compiler and timing faults and their fixes
(\src{futhark-rust-backend/README.md}; \src{docs/reports/2026-10-29_futhark_reruns.md};
\src{docs/PLAN_futhark_rust_wgpu.md}, section 4; fuller commit \texttt{3481e128}).}
\label{tab:fixes}
\small
\begin{tabularx}{\columnwidth}{@{}YY@{}}
\toprule
fault & fix and measurement\\
\midrule
SetMem of device or shared memory in a kernel loop or conditional (143 corpus programs) & patch 0008: a memory variable that can name one binding is renamed to it; otherwise a selector picks among bindings. 10 of 10 SetMem test programs identical to \texttt{futhark c}\\
a unit scalar as a kernel argument (192 programs) & patch 0007\\
flatten pairs every SegHist operator with the first one's shapes (38 programs; also in Futhark 0.28.1) & patch 0009; reported upstream with a test case\\
all compile failures & 376 of 376 corpus programs compile on the webgpu and rust paths; the default compiler is the fork with all 9 patches\\
timed runs re-uploaded the 126\,MB input; the GPU clock stepped down in the gaps & one upload per timed command, runs interleaved one by one. A/A gate: 7/22 batch, 12/22 rust, 9/22 single before; 22/22 on all three after\\
timing work bound counted loop trips, not nodes; Fable's exact program left untimed & bound $5.9\times10^{10}$ node executions\\
generation 0 differed by path & the single path's unpatched compiler wrote no C for 4 programs (SetMem); same compiler on every path\\
\bottomrule
\end{tabularx}
\end{table}

%======================================================================
\section{Discussion}
\paragraph{The fitness was the problem.} The first study's result, evolution
no better than random search, was not a property of the kingdom or of GEP.
The fitness had no slope towards the target, and on both hard benchmarks it
sloped the wrong way. The study's controls (A/A gates, planted slowdowns, the
sign test) protected the speed claims, but nothing tested the quality landscape
until the review. The slope check is now the first gate, before any run. Once
the fitness passed it, evolution's best and median fell in every run, and
random search's median did not.

\paragraph{Small targets do not test search.} On two-node targets random search
wins, because a random gene expresses the target often. A fair test of H2 needs
targets that random search cannot hit by chance. Histogram and pathfinder are
such targets, and neither regime has yet found them exactly.

\paragraph{The e-graph is a cost operator.} Every rule in every family is an
identity. In stage 1 its children always have their parent's fitness. On A7 it
produced the only cost improvement, in generation 0. On reduce, scan and
histogram it found no variants of the published programs, and on pathfinder
it found 23.

\paragraph{Where the speed came from.} Two programs faster than a published
one have been measured: the kingdom's own pathfinder (5\,\%) and Fable's
pathfinder (4--6\,\%). Neither came from evolution. The matrix family's win is
a kernel the OpenCL backend could build, found with the e-graph's help; on our
rust path the reference also runs, so that win depends on the backend.

%======================================================================
\section{Limitations and open work}
\begin{itemize}
\item \textbf{No evolved Futhark program is faster than its starting point or
  than the published program.} No rerun found an exact program, so no rerun
  reached stage 2.
\item \textbf{Few seeds.} The reruns used 2 seeds per benchmark, and two
  histogram R runs stopped at the time bound before 1,400 programs; they are
  compared at the shorter run's count. The study used 3 seeds where H1 asked
  for 5.
\item \textbf{The slope check's bar was changed after measurements were seen}
  (99\,\% to 95\,\%), with the reason recorded. The chosen formula's histogram
  Spearman is small ($+0.079$).
\item \textbf{Timing fault 13 is only partly fixed.} The verdict still compares a
  median with the least of 22 A/A ratios, and pathfinder is still timed on a
  quality case (\texttt{timing\_n = 0}, \src{src/futhark_sig.rs:163}).
\item \textbf{Operators.} Inversion and transposition have rates but no runs.
  Elite edits by a language model are planned, not built.
\item \textbf{One machine and one GPU family}: an Apple M3 Max under Metal.
  Vulkan was not measured.
\item \textbf{The matrix family is one kernel} (A7). The extension to vector
  classes (softmax, layer normalisation) waits for a second win, and an
  \texttt{exp} row would break the bit-exact contract.
\item \textbf{The signature family is integer-only.} Of 1,575 functions in the
  221 benchmark files, 10 read into the kingdom; floats, tuples and records are
  outside it.
\item \textbf{The pixel family produced no picture from random}, and is parked.
\item \textbf{Documentation lag.} Some documents disagree with the code. The
  runner's header still lists the study's seven quality inputs, the matrix
  README says 27 e-graph rules where the table has 54, and the phylu short paper
  describes the two-band fitness. This paper follows the code.
\end{itemize}

%======================================================================
\section{Conclusion}
The Futhark kingdom shows that a production array language can be made a
multi-typed GEP kingdom in which every gene compiles. Shapes, index ranges,
scopes and work bounds can all be types, generated offline and checked. The
evaluation path now scores the code that would ship, on the device, and the three
paths agree bit for bit. The search itself has not yet produced a better
Futhark program. Our first comparison found evolution no better than random
search, and the cause was a fitness with no slope towards the target. With that
fitness corrected and checked before each run, evolution beat random search in
4 of 4 pairs. The next tests are exact discovery on histogram and pathfinder,
and an improvement run that is gated by correctness and timed on the corrected
paths.

\begin{acks}
The design of multi-typed GEP, the geneframe, HFF and TrueNorth, the staged
fitness, the ring canvas and the correctness gate are Andrew Morgan's. The review
of the 22 faults was an independent reading of the code and ledgers.
\end{acks}

\bibliographystyle{ACM-Reference-Format}
\bibliography{futhark_multityped_gep}

\end{document}
